% !PRINTED 424
\weberpage{424}

%<<CONT
it follows by (1) and \S.\,111, (12):
\weq{7}{Y = \lambda v + \mu vu,}
and if, as in \S.\,111, we set
\weq{8}{\frac{H^3}{f^5} = z,\quad \frac{T^2}{f^5} = z_1,\quad z + z_1 = 1728}
then the expressions (3), (4), (5) pass into the following:
\weq{9}{az = 8\lambda^3 + \lambda^2\mu + \frac{72\lambda\mu^2 + \mu^3}{z_1}}
\[
bz = -\lambda^4 + \frac{18\lambda^2\mu^2 + \lambda\mu^3}{z_1} + 27\frac{\mu^4}{z_1^2}
\]
\[
cz = \lambda^5 - 10\frac{\lambda^3\mu^2}{z_1} + \frac{45\lambda\mu^4 + \mu^5}{z_1^2},
\]
if $Y$ is the root of the equation
\weq{10}{Y^5 + 5aY^2 + 5bY + c = 0}
This equation has the form of a principal equation of degree 5, and it can represent every principal equation if we may take $a, b, c$ arbitrarily. If, then, $a, b, c$ are assumed to be arbitrarily given, we have to determine from (9) the three quantities $\lambda, \mu, z_1$ (and $z = 1728 - z_1$), and it is a very remarkable peculiarity of this system of equations that these unknowns can be expressed \gesperrt{rationally} in terms of the given quantities $a, b, c$ and the discriminant of the equation (10). To prove this, we first derive simpler equations from (9).

If we multiply the second of them by $\lambda$ and add it to the third, it follows that
\weq{11}{z_1(\lambda b + c) = \mu^2a,}
and if we multiply the third by $\lambda$ and the second by $\mu^2 : z_1$ and subtract,
\weq{12}{z\left(\lambda c - \frac{\mu^2}{z_1}b\right) = \left(\lambda^2 - 3\frac{\mu^2}{z_1}\right)^3.}

Likewise, from the first and second we obtain
\[
z\,\frac{a\lambda + 8b}{\mu} = \lambda^3 + \frac{216\lambda^2\mu}{z_1} + 9\frac{\lambda\mu^2}{z_1} + \frac{216\mu^3}{z_1^2}.
\]
